% !TeX root = book.tex
% Nội dung lời nói đầu

Sự phát triển nhanh chóng của thương mại điện tử và logistics thông minh đã làm gia tăng nhu cầu về các giải pháp giao hàng chặng cuối hiệu quả và có tính thông minh. Giao hàng bằng máy bay không người lái đã trở thành một phương án đầy triển vọng để vận chuyển các kiện hàng nhẹ trong những môi trường như khuôn viên thông minh, khu công nghiệp, khu dân cư và vùng sâu vùng xa.\\
Dự án Nền tảng quản lý giao hàng bằng drone tích hợp AI được đề xuất dưới dạng một nền tảng quản lý giao hàng bằng máy bay không người lái có tích hợp trí tuệ nhân tạo. Hệ thống tập trung vào việc quản lý hoạt động giao hàng thay vì trực tiếp điều khiển máy bay không người lái. Nền tảng được thiết kế để hỗ trợ toàn bộ vòng đời giao hàng, bao gồm quản lý yêu cầu giao hàng, quản lý kiện hàng, thực hiện giao hàng, theo dõi giao hàng, quản lý trạm hạ cánh và xác nhận giao hàng. Bên cạnh đó, các dịch vụ hỗ trợ bởi trí tuệ nhân tạo được bổ sung nhằm dự kiến thời gian giao hàng, tóm tắt quá trình giao hàng, hỗ trợ khách hàng và phân tích hoạt động.\\
Là một thành phần quan trọng của toàn bộ hệ thống, cơ sở dữ liệu có nhiệm vụ tổ chức và duy trì dữ liệu phát sinh trong các quy trình nghiệp vụ. Một cơ sở dữ liệu được thiết kế tốt là điều kiện cần thiết để bảo đảm tính nhất quán, tính toàn vẹn, khả năng truy vết, tính bảo mật và khả năng truy cập hiệu quả cho tầng ứng dụng.\\
Tài liệu này tập trung vào thiết kế cơ sở dữ liệu cho nền tảng Nền tảng quản lý giao hàng bằng drone tích hợp AI. Phạm vi trình bày bao gồm phân tích yêu cầu dữ liệu, xác định các thực thể và mối quan hệ, thiết kế cơ sở dữ liệu ở mức khái niệm và mức logic, thiết kế lược đồ quan hệ, các ràng buộc toàn vẹn, lập chỉ mục và triển khai cơ sở dữ liệu bằng PostgreSQL.\\
